% 第 23 章：全称类型
% Source: source/split/chapters/ch23-universal-types_p321-p343.pdf
% Original print pages: 339--362
% Status: 待独立初审

\chapter{全称类型}
\label{chap:23}

上一章考察了 ML 中受限的 let 多态。本章转向一种更一般的参数多态系统：
\termfirst{System F}{System F}，也称\termfirst{多态 lambda 演算}{polymorphic
lambda-calculus}或二阶 lambda 演算。

\section{动机}
\label{sec:23.1}

在简单类型 lambda 演算中，可以写出无穷多个“加倍”函数：
\[
\begin{array}{l}
\taplterm{doubleNat}
 =\lambda f:\tapltype{Nat}\to\tapltype{Nat}.\,
   \lambda x:\tapltype{Nat}.\,f(f\,x),\\
\taplterm{doubleRcd}
 =\lambda f:\{l:\tapltype{Bool}\}\to\{l:\tapltype{Bool}\}.\,
   \lambda x:\{l:\tapltype{Bool}\}.\,f(f\,x),\\
\taplterm{doubleFun}
 =\lambda f:(\tapltype{Nat}\to\tapltype{Nat})\to
              (\tapltype{Nat}\to\tapltype{Nat}).\,
   \lambda x:\tapltype{Nat}\to\tapltype{Nat}.\,f(f\,x)
\end{array}
\]
三者接受不同类型的参数，行为却完全相同；除类型标注外，程序文本也相同。
若同一程序需要在不同类型上使用这一操作，简单类型系统迫使我们复制定义。

\begin{keystatement}
\textbf{抽象原则：}程序中每项重要功能都应只在源码的一处实现。若多段代码执行
相似功能，通常应抽取其中变化的部分，把它们合并为一个定义。
\end{keystatement}

这里变化的正是类型。因此，我们需要从项中抽象出类型，并在使用该项时再提供
具体类型。

\footnotetext{本章主要研究纯 System F（\cref{fig:23-1}）；示例会自由使用
此前介绍的扩展。原书相应的 OCaml 实现为 \texttt{fullpoly}；涉及序对和列表
的部分示例需要 \texttt{fullomega}。}

\section{多态的几种形式}
\label{sec:23.2}

允许同一段代码以多种类型使用的系统统称为\termfirst{多态系统}{polymorphic
system}。现代语言中常见的多态形式包括以下几种。

\begin{description}
  \item[参数多态（parametric polymorphism）]
    以类型变量代替具体类型编写通用代码，并在需要时以具体类型实例化。
    参数多态定义具有一致的行为：不同类型实例执行同一种算法。
  \item[特设多态（ad-hoc polymorphism）]
    同一个名称在不同类型下可以采用不同实现。最常见的例子是
    \termfirst{重载}{overloading}：编译器或运行时系统根据参数类型选择具体
    实现。多方法分派、受限的运行时类型分析等机制也属于这一大类。
  \item[子类型多态（subtype polymorphism）]
    借助 T-Sub，同一个项可以获得多个类型；把它看作某个超类型时，会有选择地
    忘掉关于该项行为的部分信息。
\end{description}

本章研究参数多态中能力最强的一种形式：非直谓多态，也称一等多态。实践中
更常见的是 ML 风格的 let 多态：它把多态限制在 let 绑定处，不允许函数接收
多态值作为参数，换来较自然的自动类型重构（\chapref{22}）。

上述类别并不互斥。一门语言可以同时具有多种多态。例如 Standard ML 兼有参数
多态和少量内建运算重载；Java 则长期以子类型、重载和受限的运行时类型测试为
主，后来又加入泛型。不同社群单说“多态”时含义也可能不同：函数式编程语境
通常指参数多态，面向对象语境则常指子类型多态，而以“泛型”指参数多态。

\section{System F}
\label{sec:23.3}

System F 最早由 Jean-Yves Girard（1972）在证明论研究中提出；John Reynolds
（1974）随后独立提出了能力基本相同的多态 lambda 演算。借助 Curry--Howard
对应，它对应于二阶直觉主义逻辑：量化对象不再限于项，也可以包括类型。

System F 在简单类型 lambda 演算上增加两种项：
\[
\lambda X.\,t
\qquad\text{和}\qquad
t[T]
\]
前者是\termfirst{类型抽象}{type abstraction}，从项中抽象出类型变量 \(X\)；
后者是\termfirst{类型应用}{type application}或实例化，把具体类型 \(T\)
提供给类型抽象。二者相遇时发生类型层的 beta 归约：
\[
(\lambda X.\,t_{11})[T_2]\longrightarrow[X\mapsto T_2]t_{11}
\]

例如，多态恒等函数
\[
\taplterm{id}=\lambda X.\,\lambda x:X.\,x
\]
应用于 \(\tapltype{Nat}\) 后得到
\(\lambda x:\tapltype{Nat}.\,x\)。

用全称类型
\[
\forall X.\,T
\]
给类型抽象分类。恒等函数的类型是
\(\forall X.\,X\to X\)：无论给它什么类型 \(T\)，实例
\(\taplterm{id}[T]\) 都具有类型 \(T\to T\)。

\begin{figure}[!ht]
\centering
\begin{minipage}{0.96\linewidth}
\textbf{类型}
\[
T ::= X \mid T\to T \mid \forall X.\,T
\]
\textbf{项与值}
\[
\begin{array}{rcl}
t&::=&x\mid\lambda x:T.\,t\mid t\,t\mid\lambda X.\,t\mid t[T],\\
v&::=&\lambda x:T.\,t\mid\lambda X.\,t
\end{array}
\]
\textbf{求值}
\[
\inferrule*[right=E-TApp]
 {t_1\longrightarrow t_1'}
 {t_1[T_2]\longrightarrow t_1'[T_2]}
\qquad
\inferrule*[right=E-TAppTabs]{ }
 {(\lambda X.\,t_{11})[T_2]\longrightarrow[X\mapsto T_2]t_{11}}
\]
\textbf{赋型}
\[
\inferrule*[right=T-TAbs]
 {\Gamma,X\vdash t_2:T_2}
 {\Gamma\vdash\lambda X.\,t_2:\forall X.\,T_2}
\qquad
\inferrule*[right=T-TApp]
 {\Gamma\vdash t_1:\forall X.\,T_{11}}
 {\Gamma\vdash t_1[T_2]:[X\mapsto T_2]T_{11}}
\]
\end{minipage}
\caption{多态 lambda 演算（System F）}
\label{fig:23-1}
\end{figure}

T-TAbs 的子推导把 \(X\) 加入上下文，用来记录其作用域。沿用
\cref{sec:5.3}的约定，新绑定的项变量和类型变量均不与上下文中已有名称
重复；必要时可以一致地重命名绑定变量。纯演算省略了记录、自然数、布尔值、
let 和 fix 等扩展，后文示例仍会自由使用它们。

\section{示例}
\label{sec:23.4}

\subsection*{热身}

多态恒等函数和加倍函数为
\[
\begin{array}{rcl}
\taplterm{id}&=&\lambda X.\,\lambda x:X.\,x,\\
\taplterm{double}&=&
\lambda X.\,\lambda f:X\to X.\,\lambda a:X.\,f(f\,a)
\end{array}
\]
类型分别是
\[
\forall X.\,X\to X
\qquad\text{和}\qquad
\forall X.\,(X\to X)\to X\to X.
\]
先对类型参数实例化，再对普通参数应用：
\[
\taplterm{double}[\tapltype{Nat}]
  (\lambda x:\tapltype{Nat}.\,\taplterm{succ}(\taplterm{succ}\,x))\ 3
\longrightarrow^*7.
\]

System F 还允许多态自应用。令
\[
\taplterm{selfApp}
=\lambda x:\forall X.\,X\to X.\,
   x[\forall X.\,X\to X]\ x
\]
则
\[
\taplterm{selfApp}:
(\forall X.\,X\to X)\to(\forall X.\,X\to X).
\]
这里先把 \(x\) 的全称类型实例化为它自身的类型，所得函数便可以接收原来的
\(x\)。同理，
\[
\taplterm{quadruple}
=\lambda X.\,
  \taplterm{double}[X\to X]\bigl(\taplterm{double}[X]\bigr)
\]
把加倍操作应用于自身，得到四次应用。

\begin{exercise}[\(\star\star\)]
\label{ex:23.4.1}
使用\cref{fig:23-1}的规则，验证上述各项确实具有所标出的类型。
\end{exercise}
\taplsolutionlink{sol:translator:23.4.1}{23.4.1}

\subsection*{多态列表}

若语言提供列表构造子，常用操作可具有如下全称类型：
\[
\begin{array}{rcl}
\taplterm{nil}&:&\forall X.\,\tapltype{List}\,X,\\
\taplterm{cons}&:&\forall X.\,X\to\tapltype{List}\,X\to\tapltype{List}\,X,\\
\taplterm{isnil}&:&\forall X.\,\tapltype{List}\,X\to\tapltype{Bool},\\
\taplterm{head}&:&\forall X.\,\tapltype{List}\,X\to X,\\
\taplterm{tail}&:&\forall X.\,\tapltype{List}\,X\to\tapltype{List}\,X
\end{array}
\]
这与\cref{sec:11.12}为列表专门编写的规则表达同一约束，但列表现在可以被
看作提供若干多态常量的库，而不必烘焙进核心语言。

\[
\begin{aligned}
\taplterm{map}={}&\lambda X.\,\lambda Y.\,
 \lambda f:X\to Y.\,
 \taplterm{fix}\bigl(
   \lambda m:\tapltype{List}\,X\to\tapltype{List}\,Y.\,
   \lambda l:\tapltype{List}\,X.\\[-1mm]
 &\quad\textbf{if }\taplterm{isnil}[X]\ l
   \textbf{ then }\taplterm{nil}[Y]\\
 &\quad\textbf{ else }\taplterm{cons}[Y]\,
      (f(\taplterm{head}[X]\ l))\,
      (m(\taplterm{tail}[X]\ l))\bigr)
\end{aligned}
\]
它具有类型
\[
\forall X.\,\forall Y.\,(X\to Y)\to\tapltype{List}\,X\to
\tapltype{List}\,Y.
\]

\begin{exercise}[\(\star\star\)]
\label{ex:23.4.2}
验证 \(\taplterm{map}\) 具有上述类型。
\end{exercise}
\taplsolutionlink{sol:translator:23.4.2}{23.4.2}

\begin{exercise}[推荐，\(\star\star\)]
\label{ex:23.4.3}
仿照 \(\taplterm{map}\)，写出多态列表反转函数
\[
\taplterm{reverse}:\forall X.\,\tapltype{List}\,X\to\tapltype{List}\,X.
\]
\end{exercise}
\taplauthorsolutionlink{sol:author:23.4.3}{23.4.3}

\begin{exercise}[\(\star\star\)]
\label{ex:23.4.4}
写出简单的多态排序函数
\[
\taplterm{sort}:\forall X.\,(X\to X\to\tapltype{Bool})\to
\tapltype{List}\,X\to\tapltype{List}\,X.
\]
第一个参数是 \(X\) 元素的比较函数。
\end{exercise}
\taplsolutionlink{sol:translator:23.4.4}{23.4.4}

\subsection*{Church 编码}

\cref{sec:5.2}曾在无类型 lambda 演算中用函数编码布尔值、自然数和列表。
System F 也能给这些编码加上类型。Church 布尔值的类型为
\[
\tapltype{CBool}=\forall X.\,X\to X\to X
\]
两个值为
\[
\begin{array}{rcl}
\taplterm{tru}&=&\lambda X.\,\lambda t:X.\,\lambda f:X.\,t,\\
\taplterm{fls}&=&\lambda X.\,\lambda t:X.\,\lambda f:X.\,f.
\end{array}
\]
例如
\[
\taplterm{not}
=\lambda b:\tapltype{CBool}.\,
 \lambda X.\,\lambda t:X.\,\lambda f:X.\,b[X]\ f\ t.
\]

\begin{exercise}[推荐，\(\star\)]
\label{ex:23.4.5}
写出接受两个 \(\tapltype{CBool}\) 参数并计算合取的
\(\taplterm{and}\)。
\end{exercise}
\taplauthorsolutionlink{sol:author:23.4.5}{23.4.5}

Church 自然数的类型为
\[
\tapltype{CNat}=\forall X.\,(X\to X)\to X\to X.
\]
例如
\[
\begin{array}{rcl}
\taplterm{c0}&=&\lambda X.\,\lambda s:X\to X.\,\lambda z:X.\,z,\\
\taplterm{c1}&=&\lambda X.\,\lambda s:X\to X.\,\lambda z:X.\,s\,z,\\
\taplterm{c2}&=&\lambda X.\,\lambda s:X\to X.\,\lambda z:X.\,s(s\,z).
\end{array}
\]
后继和加法可以定义为
\[
\begin{aligned}
\taplterm{csucc}
  &={}\lambda n:\tapltype{CNat}.\,\lambda X.\,
       \lambda s:X\to X.\,\lambda z:X.\,s(n[X]\ s\ z),\\
\taplterm{cplus}
  &={}\lambda m:\tapltype{CNat}.\,\lambda n:\tapltype{CNat}.\,
       \lambda X.\,\lambda s:X\to X.\,\lambda z:X.\,
       m[X]\ s\,(n[X]\ s\ z).
\end{aligned}
\]

\begin{exercise}[推荐，\(\star\star\)]
\label{ex:23.4.6}
写出函数 \(\taplterm{iszero}\)，使其应用于 \(\taplterm{c0}\) 时返回
\(\taplterm{tru}\)，应用于其他 Church 数时返回 \(\taplterm{fls}\)。
\end{exercise}
\taplauthorsolutionlink{sol:author:23.4.6}{23.4.6}

\begin{exercise}[\(\star\star\)]
\label{ex:23.4.7}
验证
\[
\begin{aligned}
\taplterm{ctimes}
 &=\lambda m:\tapltype{CNat}.\,\lambda n:\tapltype{CNat}.\,
   \lambda X.\,\lambda s:X\to X.\,n[X](m[X]s),\\
\taplterm{cexp}
 &=\lambda m:\tapltype{CNat}.\,\lambda n:\tapltype{CNat}.\,
   \lambda X.\,n[X\to X](m[X])
\end{aligned}
\]
均具有类型
\(\tapltype{CNat}\to\tapltype{CNat}\to\tapltype{CNat}\)，并说明它们
为何分别表示乘法和幂运算。
\end{exercise}
\taplsolutionlink{sol:translator:23.4.7}{23.4.7}

\begin{exercise}[推荐，\(\star\star\)]
\label{ex:23.4.8}
证明
\[
\tapltype{PairNat}
=\forall X.\,(\tapltype{CNat}\to\tapltype{CNat}\to X)\to X
\]
可以表示自然数对，并定义构造函数以及第一、第二投影。
\end{exercise}
\taplauthorsolutionlink{sol:author:23.4.8}{23.4.8}

\begin{exercise}[推荐，\(\star\star\star\)]
\label{ex:23.4.9}
利用\cref{ex:23.4.8}的数对写出 Church 数的前驱函数。提示：定义
\(f(i,j)=(i+1,i)\)，把它从 \((0,0)\) 开始迭代 \(n\) 次，再取第二分量。
\end{exercise}
\taplauthorsolutionlink{sol:author:23.4.9}{23.4.9}

\begin{exercise}[\(\star\star\star\)]
\label{ex:23.4.10}
设 \(k=\lambda x.\lambda y.x\)、\(i=\lambda x.x\)。为 Velmans 的无类型
前驱项
\[
\lambda n.\lambda s.\lambda z.\,
n(\lambda p.\lambda q.\,q(p\,s))(k\,z)i
\]
补上类型抽象、类型应用和参数类型，使其成为 System F 项；并说明其工作原理。
\end{exercise}
\taplauthorsolutionlink{sol:author:23.4.10}{23.4.10}

\subsection*{列表编码}

含 \(X\) 元素的 Church 列表可以看作它自己的右折叠函数：
\[
\tapltype{List}\,X=\forall R.\,(X\to R\to R)\to R\to R.
\]
于是
\[
\begin{aligned}
\taplterm{nil}
 &={}\lambda X.\,\lambda R.\,\lambda c:X\to R\to R.\,\lambda n:R.\,n,\\
\taplterm{cons}
 &={}\lambda X.\,\lambda hd:X.\,\lambda tl:\tapltype{List}\,X.\,
   \lambda R.\,\lambda c:X\to R\to R.\,\lambda n:R.\,
   c\ hd\ (tl[R]\ c\ n),\\
\taplterm{isnil}
 &={}\lambda X.\,\lambda l:\tapltype{List}\,X.\,
   l[\tapltype{Bool}](\lambda hd:X.\lambda tl:\tapltype{Bool}.
   \taplterm{false})\ \taplterm{true}.
\end{aligned}
\]

为处理空列表的 \(\taplterm{head}\)，若语言带 \(\taplterm{fix}\)，可先定义
\[
\taplterm{diverge}
=\lambda X.\,\lambda\_:\tapltype{Unit}.\,
 \taplterm{fix}(\lambda x:X.\,x)
:\forall X.\,\tapltype{Unit}\to X.
\]
直接把 \(\taplterm{diverge}[X]\ \taplterm{unit}\) 作为折叠基值会在非空
列表上也过早发散。应把计算推迟到最后：
\[
\begin{aligned}
\taplterm{head}={}&
\lambda X.\,\lambda l:\tapltype{List}\,X.\\
&\bigl(l[\tapltype{Unit}\to X]\,
(\lambda hd:X.\,\lambda tl:\tapltype{Unit}\to X.\,
 \lambda\_:\tapltype{Unit}.\,hd)\,
(\taplterm{diverge}[X])\bigr)\ \taplterm{unit}.
\end{aligned}
\]
空列表产生发散函数；非空列表产生忽略 \(\taplterm{unit}\) 并返回首元素的
函数，最后一次应用才决定结果。

\begin{exercise}[\(\star\star\)]
\label{ex:23.4.11}
纯 System F 没有 \(\taplterm{fix}\)。写出另一版 \(\taplterm{head}\)，
额外接收一个参数，在列表为空时返回该参数。
\end{exercise}
\taplauthorsolutionlink{sol:author:23.4.11}{23.4.11}

\begin{exercise}[推荐，\(\star\star\star\)]
\label{ex:23.4.12}
在不含 \(\taplterm{fix}\) 的纯 System F 中，定义
\[
\taplterm{insert}:
\forall X.\,(X\to X\to\tapltype{Bool})\to
\tapltype{List}\,X\to X\to\tapltype{List}\,X,
\]
并用它构造列表排序函数。
\end{exercise}
\taplauthorsolutionlink{sol:author:23.4.12}{23.4.12}

\section{基本性质}
\label{sec:23.5}

System F 的类型安全性质与简单类型 lambda 演算基本相同。

\begin{theorem}[保持性]
\label{thm:23.5.1}
若 \(\Gamma\vdash t:T\) 且 \(t\longrightarrow t'\)，则
\(\Gamma\vdash t':T\)。
\end{theorem}
\taplauthorsolutionlink{sol:author:23.5.1}{23.5.1}

\begin{theorem}[进展性]
\label{thm:23.5.2}
若 \(t\) 是闭的良类型项，则 \(t\) 或为值，或存在 \(t'\) 使
\(t\longrightarrow t'\)。
\end{theorem}
\taplauthorsolutionlink{sol:author:23.5.2}{23.5.2}

System F 还和简单类型 lambda 演算一样具有规范化性质。这一结论并不显然：
即使不借助 \(\taplterm{fix}\)，纯 System F 也能表达排序等相当复杂的计算。
Girard 在其博士论文中以\chapref{12}方法的推广证明了这一性质。

\begin{theorem}[规范化]
\label{thm:23.5.3}
每个良类型 System F 项都可规范化。
\end{theorem}

\section{擦除、可赋型性与类型重构}
\label{sec:23.6}

System F 项可以通过删去类型标注、类型抽象和类型应用而映射到无类型项：
\[
\begin{array}{rcl}
\mathsf{erase}(x)&=&x,\\
\mathsf{erase}(\lambda x:T_1.\,t_2)&=&\lambda x.\mathsf{erase}(t_2),\\
\mathsf{erase}(t_1\,t_2)&=&\mathsf{erase}(t_1)\,\mathsf{erase}(t_2),\\
\mathsf{erase}(\lambda X.\,t_2)&=&\mathsf{erase}(t_2),\\
\mathsf{erase}(t_1[T_2])&=&\mathsf{erase}(t_1).
\end{array}
\]
若存在良类型 System F 项 \(t\) 使 \(\mathsf{erase}(t)=m\)，则称无类型项 \(m\)
在 System F 中可赋型。类型重构问题要求由 \(m\) 找出
这样的 \(t\)，或判断不存在。

\begin{theorem}[Wells，1994]
\label{thm:23.6.1}
给定闭的无类型 lambda 项 \(m\)，判断是否存在良类型 System F 项 \(t\)
满足 \(\mathsf{erase}(t)=m\)，是不可判定的。
\end{theorem}

即使保留所有普通类型标注以及类型抽象的位置，只抹掉类型应用的实参，这种
部分重构仍不可判定。令
\[
\begin{array}{rcl}
\mathsf{erase}_p(x)&=&x,\\
\mathsf{erase}_p(\lambda x:T_1.\,t_2)&=&\lambda x:T_1.\mathsf{erase}_p(t_2),\\
\mathsf{erase}_p(t_1t_2)&=&\mathsf{erase}_p(t_1)\mathsf{erase}_p(t_2),\\
\mathsf{erase}_p(\lambda X.\,t_2)&=&\lambda X.\mathsf{erase}_p(t_2),\\
\mathsf{erase}_p(t_1[T_2])&=&\mathsf{erase}_p(t_1)[\,].
\end{array}
\]

\begin{theorem}[Boehm，1985，1989]
\label{thm:23.6.2}
给定类型应用位置已标出但类型实参被省略的闭项 \(s\)，判断是否存在良类型
System F 项 \(t\) 使 \(\mathsf{erase}_p(t)=s\)，是不可判定的。
\end{theorem}

\translatornote{作者勘误修正了本节部分擦除函数的定义。这里采用修正后的
版本：普通类型标注和类型抽象均保留，只有类型应用的实参被替换为空方括号。
对于另一些介于完整标注与完全擦除之间的精确重构问题，已知结果更细致，不能
由本节两条定理一概而论。}

实际语言通常采用受限的局部类型推断、显式标注与推断相结合，或把可推断的
多态限制在较温和的片段。它们不恢复完整 System F 类型重构的可判定性，而是
选择足够实用且行为可预测的子问题。

\begin{exercise}[\(\star\star\star\star\)]
\label{ex:23.6.3}
不用规范化定理，只从赋型规则出发证明
\(\Omega=(\lambda x.x\,x)(\lambda y.y\,y)\) 在 System F 中不可赋型。
题目要求先把任意良类型装饰整理为外层类型抽象/应用与“暴露项”的组合，
再分析 \(x\,x\) 对类型结构造成的矛盾。
\end{exercise}
\taplauthorsolutionlink{sol:author:23.6.3}{23.6.3}

\section{擦除与求值次序}
\label{sec:23.7}

\cref{fig:23-1}采用\termfirst{类型传递语义}{type-passing semantics}：
运行时会把类型实参替换进类型抽象体。实际编译器通常希望类型检查后擦除类型。
不过，在含副作用的按值调用语言中，简单擦除会改变求值次序。例如
\[
\texttt{let }f=(\lambda X.\,\taplterm{error})\texttt{ in }0
\]
求值为 \(0\)，因为类型抽象是值，其内部错误不会立即执行；若直接擦除类型
抽象，就得到 \(\texttt{let }f=\taplterm{error}\texttt{ in }0\)，会立刻
抛出错误。

适合按值调用的擦除应把类型抽象保留为无类型的零信息函数，把类型应用变成向它
传递任意哑值：
\[
\begin{array}{rcl}
\mathsf{erase}_v(x)&=&x,\\
\mathsf{erase}_v(\lambda x:T_1.\,t_2)&=&\lambda x.\mathsf{erase}_v(t_2),\\
\mathsf{erase}_v(t_1t_2)&=&\mathsf{erase}_v(t_1)\mathsf{erase}_v(t_2),\\
\mathsf{erase}_v(\lambda X.\,t_2)&=&\lambda\_.\mathsf{erase}_v(t_2),\\
\mathsf{erase}_v(t_1[T_2])&=&\mathsf{erase}_v(t_1)\,\taplterm{dummyv}.
\end{array}
\]
其中 \(\taplterm{dummyv}\) 可以取任意无类型值，如
\(\taplterm{unit}\)。这样，类型抽象原来阻止函数体立即求值的作用仍然保留。

\begin{exercise}[\(\star\star\)]
\label{ex:23.7.1}
把原书第 335 页关于引用与不安全 let 多态的反例，写成带引用的 System F
项。
\end{exercise}
\taplauthorsolutionlink{sol:author:23.7.1}{23.7.1}

\begin{theorem}
\label{thm:23.7.2}
若 \(\mathsf{erase}_v(t)=u\)，则有且仅有以下两种情况之一：
\begin{enumerate}
  \item \(t\) 与 \(u\) 都是各自求值关系下的范式；
  \item \(t\longrightarrow t'\)、\(u\longrightarrow u'\)，并且
        \(\mathsf{erase}_v(t')=u'\)。
\end{enumerate}
\end{theorem}

\section{System F 的受限片段}
\label{sec:23.8}

System F 优雅而强大，但完整类型重构不可判定。语言设计因此常选择能力较弱、
重构更可控的片段。ML 的 let 多态也称\termfirst{前束多态}{prenex
polymorphism}：类型变量只实例化为不含量词的单型，而且量化类型不出现在箭头
左侧。

另一种常见限制是\termfirst{秩 2 多态}{rank-2 polymorphism}。把类型画成树；
若从根到任一 \(\forall\) 的路径，在两个或更多箭头处进入参数一侧，该类型就
超过秩 2。例如
\[
(\forall X.\,X\to X)\to\tapltype{Nat}
\]
属于秩 2，而
\[
((\forall X.\,X\to X)\to\tapltype{Nat})\to\tapltype{Nat}
\]
不属于。秩 2 系统能给比 ML 片段更多的无类型项赋型；其重构复杂度与 ML
相同，而秩 3 及以上的 System F 重构不可判定。

\section{参数性}
\label{sec:23.9}

观察
\[
\tapltype{CBool}=\forall X.\,X\to X\to X.
\]
这种类型的值必须先是类型抽象，再接收两个 \(X\) 型参数，最后返回一个 \(X\)
型结果。因为 \(X\) 是未知类型，函数无法自行构造 \(X\) 值，只能返回收到的
两个参数之一。因此，忽略求值上等价的写法后，该类型的值只能表现得像
\(\taplterm{tru}\) 或 \(\taplterm{fls}\)。

这一现象来自\termfirst{参数性}{parametricity}：参数多态程序不能根据未知
类型的内部结构改变行为，其类型因而蕴含很强的行为约束。Wadler（1989）把从
多态类型直接推导出的此类性质称为“免费定理”。

\section{非直谓性}
\label{sec:23.10}

若一个定义中的量词范围包含正在定义的对象本身，就称该定义是
\termfirst{非直谓的}{impredicative}。例如
\[
T=\forall X.\,X\to X
\]
中的 \(X\) 遍历所有类型，其中也包括 \(T\) 本身；所以 \(T\) 型项可以在
\(T\) 处实例化。ML 多态通常称为\termfirst{直谓的}{predicative}或分层的，
因为类型变量只遍历不含量词的单型。

“直谓”与“非直谓”来自逻辑和集合论。Russell 与 Poincaré 曾用这一划分排除
通过包含自身的量化范围定义集合所导致的恶性循环。System F 的非直谓量化并不
因此不一致；它的证明论解释和规范化性质恰好说明，该类型系统仍然具有严密的
逻辑基础。

\section{注记}
\label{sec:23.11}

关于 System F 的进一步入门材料，可参见 Reynolds（1990）及其
\emph{Theories of Programming Languages}（1998）。
